% NE 630 -- Lesson 24 handout
% Build: make L24 (XeLaTeX). Original Downloads draft is preserved.
% Set the single solution flag below true only for a private solution copy.
\documentclass[compact]{ne630handout}
\handoutsolutionsfalse
\hypersetup{pdftitle={NE 630 -- Lesson 24: Linear Reactivity Model},
  pdfauthor={Jeremy Roberts},pdfsubject={Equilibrium batch refueling and burnup}}
\newcommand{\Ufive}{{}^{235}\mathrm{U}}
\newcommand{\Ueight}{{}^{238}\mathrm{U}}
\newcommand{\Punine}{{}^{239}\mathrm{Pu}}

\begin{document}
\fontsize{8.85}{10.35}\selectfont
\RaggedRight
\setlength{\abovedisplayskip}{3.5pt}
\setlength{\belowdisplayskip}{3.5pt}
\setlength{\abovedisplayshortskip}{2.5pt}
\setlength{\belowdisplayshortskip}{2.5pt}

\Lesson{24}{Linear reactivity model}
 {Linear Reactivity Model supplement; review Lesson 23.}
 {How do batch refueling strategies change cycle length and discharge burnup?}
\textbf{Objective:} Estimate cycle lengths and discharge burnups using the linear reactivity model.

\noindent\begin{minipage}[t]{0.485\textwidth}\vspace{0pt}\RaggedRight
\Step{1}{Review: reactivity and depletion}
For constant total flux and one-group cross sections, constant $N^{\Ueight}$, and $N^{\Punine}(0)=0$, Lesson 23 gives
\begin{align*}
 N^{\Ufive}(t)&=N^{\Ufive}(0)e^{-\bar\sigma_a^{\Ufive}\phi t},\\
 N^{\Ueight}(t)&=N^{\Ueight}(0),\\
 N^{\Punine}(t)&=
 \frac{\bar\sigma_\gamma^{\Ueight}}{\bar\sigma_a^{\Punine}}
 N^{\Ueight}(0)
 \left[1-e^{-\bar\sigma_a^{\Punine}\phi t}\right].
\end{align*}
Neglecting control and fission-product absorption,
\begin{equation*}
 k(t)=
 \frac{\left[\bar\nu_T^{\Ufive}\bar\sigma_{fT}^{\Ufive}N^{\Ufive}(t)
 +\bar\nu_T^{\Punine}\bar\sigma_{fT}^{\Punine}N^{\Punine}(t)\right]
 \varepsilon pP_{\mathrm{NL}}}
 {\bar\sigma_{aT}^{\Ufive}N^{\Ufive}(t)
 +\bar\sigma_{aT}^{\Punine}N^{\Punine}(t)
 +\Sigma_{aT}^{\Ueight}
 +\varsigma\dfrac{V_m}{V_f}\Sigma_{aT}^{m}},
 \tag{FNRP 10.39*}
\end{equation*}
with $\rho=(k-1)/k$. The asterisk makes thermal-group ($T$) cross sections and nuclide-specific neutron yields per fission explicit.

\Step{2}{Let's get some rough PWR numbers}
Let's do some back-of-the-envelope calculations to get rough (one- or two-significant-figure) numbers for a typical PWR. Use the rounded $V_F=10^7~\mathrm{cm^3}$ in parts 2--4; burnup is per initial uranium mass.

\textbf{1. Fuel volume.} A $3000~\mathrm{MW_{th}}$ PWR has about $50{,}000$ rods, each roughly $3~\mathrm{m}$ high. With $r\simeq0.5~\mathrm{cm}$,
\[
 V_F\simeq(5\times10^4)\pi r^2h
 \simeq \RevealMath[0.75in]{10^7\ \mathrm{cm^3}}.
\]

\textbf{2. Flux.} If $\bar\Sigma_f=0.05~\mathrm{cm^{-1}}$ and $\kappa\simeq3.2\times10^{-11}~\mathrm{J/fission}$, then
\[
 \phi=\frac{P}{V_F\bar\Sigma_f\kappa}
 \simeq \RevealMath[1.05in]{2\times10^{14}\ \mathrm{cm^{-2}s^{-1}}}.
\]

\textbf{3. Fuel mass.} With $\rho_{\mathrm{UO_2}}=10~\mathrm{g/cm^3}$,
\begin{align*}
 m_{\mathrm{UO_2}}&=\rho_{\mathrm{UO_2}}V_F
 \simeq\RevealMath[0.72in]{1.0\times10^5\ \mathrm{kg}},\\
 m_{\mathrm U}&\simeq\frac{238}{270}m_{\mathrm{UO_2}}
 \simeq\RevealMath[0.72in]{8.8\times10^4\ \mathrm{kg}}.
\end{align*}

\textbf{4. Burnup.} For five full-power years (not calendar years),
\begin{align*}
 B_{1\,\mathrm d}&\simeq\frac{3000~\mathrm{MWd}}{8.8\times10^4~\mathrm{kg}}
 \simeq\RevealMath[0.65in]{0.034\ \mathrm{MWd/kg}},\\
 B_{5\,\mathrm{yr}}&\simeq5(365)B_{1\,\mathrm d}
 \simeq\RevealMath[0.65in]{62\ \mathrm{MWd/kg}}.
\end{align*}
\end{minipage}
\hfill
\begin{minipage}[t]{0.485\textwidth}\vspace{0pt}\RaggedRight
\Step{3}{Use a more realistic reactivity curve}
Use the fixed-power curve from \texttt{lesson\_24.ipynb}: 4\% enriched UO$_2$ at $10.4~\mathrm{g/cm^3}$, $300~\mathrm{W/cm^3}$ fuel, and the expanded chain with Xe/Sm. The prescribed-spectrum rate ratio uses $P_{\mathrm{NL}}=1$; leakage is treated separately below.

\begin{center}
\begin{tikzpicture}
\begin{axis}[
  width=0.99\linewidth,height=2.15in,
  xmin=0,xmax=40,ymin=0.06,ymax=0.28,
  axis lines=left,
  xlabel={$B$ (MWd/kg initial HM)},ylabel={$\rho_\infty$},
  tick label style={font=\scriptsize},
  label style={font=\scriptsize},
  xtick={0,5,...,40},ytick={0.06,0.08,...,0.28},
  scaled y ticks=false,yticklabel style={/pgf/number format/fixed,/pgf/number format/precision=2},
  minor x tick num=4,minor y tick num=1,
  grid=both,major grid style={gray!45},minor grid style={gray!20}
]
% Actual fixed-power depletion history from lesson_24.ipynb, cells 2 and 5.
% Fixed spectrum (0.35,0.55,0.10); Table-3.2 XS with U238 rod shielding.
% Curve remains visible in BOTH student and instructor copies.
\addplot[KSUPurple,thick] coordinates {
  (0,0.27468133)
  (2.2726812e-05,0.27467871)
  (2.6014502e-05,0.27467832)
  (2.9777795e-05,0.27467788)
  (3.408549e-05,0.27467737)
  (3.9016342e-05,0.27467677)
  (4.4660498e-05,0.27467609)
  (5.1121145e-05,0.2746753)
  (5.8516397e-05,0.27467439)
  (6.6981457e-05,0.27467333)
  (7.6671083e-05,0.27467211)
  (8.7762424e-05,0.27467068)
  (0.00010045825,0.27466902)
  (0.00011499068,0.27466708)
  (0.00013162538,0.27466481)
  (0.00015066648,0.27466215)
  (0.00017246209,0.27465903)
  (0.00019741069,0.27465534)
  (0.00022596838,0.27465098)
  (0.00025865727,0.27464581)
  (0.00029607497,0.27463965)
  (0.00033890557,0.2746323)
  (0.00038793209,0.27462348)
  (0.00044405087,0.27461288)
  (0.00050828785,0.27460007)
  (0.00058181745,0.27458455)
  (0.00066598393,0.27456569)
  (0.00076232605,0.27454268)
  (0.00087260514,0.27451453)
  (0.00099883736,0.27448001)
  (0.0011433305,0.27443757)
  (0.0013087262,0.27438531)
  (0.0014980483,0.27432087)
  (0.0017147579,0.27424136)
  (0.0019628171,0.27414326)
  (0.0022467608,0.27402228)
  (0.0025717802,0.27387332)
  (0.0029438173,0.27369033)
  (0.0033696738,0.27346626)
  (0.0038571353,0.27319309)
  (0.0044151136,0.27286185)
  (0.0050538098,0.27246286)
  (0.0057849006,0.27198614)
  (0.006621752,0.27142193)
  (0.0075796634,0.27076163)
  (0.0086761475,0.26999885)
  (0.0099312505,0.26913083)
  (0.011367918,0.26815994)
  (0.013012417,0.2670951)
  (0.01489481,0.265953)
  (0.017049513,0.26475848)
  (0.019515918,0.26354398)
  (0.022339117,0.26234758)
  (0.025570723,0.2612096)
  (0.029269817,0.26016812)
  (0.033504028,0.25925406)
  (0.038350765,0.25848682)
  (0.043898638,0.2578716)
  (0.050249073,0.25739923)
  (0.05751817,0.25704862)
  (0.065838825,0.25679123)
  (0.075363157,0.25659635)
  (0.086265291,0.25643573)
  (0.098744542,0.25628659)
  (0.11302906,0.25613253)
  (0.12937999,0.25596284)
  (0.14809627,0.25577094)
  (0.16952008,0.25555281)
  (0.19404307,0.25530595)
  (0.2221136,0.25502884)
  (0.25424485,0.25472085)
  (0.29102425,0.25438229)
  (0.3331242,0.25401449)
  (0.38131439,0.25361979)
  (0.43647583,0.2532014)
  (0.49961701,0.25276304)
  (0.57189226,0.25230855)
  (0.65462296,0.25184122)
  (0.74932158,0.25136309)
  (0.85771943,0.25087418)
  (0.98179826,0.25037168)
  (1,0.25030318)
  (1.25,0.24941727)
  (1.5,0.248612)
  (1.75,0.24783624)
  (2,0.24706442)
  (2.25,0.24628281)
  (2.5,0.24548496)
  (2.75,0.24466781)
  (3,0.24382979)
  (3.25,0.24297118)
  (3.5,0.24209227)
  (3.75,0.24119325)
  (4,0.24027529)
  (4.25,0.23933914)
  (4.5,0.2383852)
  (4.75,0.23741466)
  (5,0.23642829)
  (5.25,0.23542649)
  (5.5,0.23441034)
  (5.75,0.23338057)
  (6,0.23233756)
  (6.25,0.23128227)
  (6.5,0.23021536)
  (6.75,0.22913722)
  (7,0.22804867)
  (7.25,0.22695033)
  (7.5,0.22584259)
  (7.75,0.22472614)
  (8,0.22360155)
  (8.25,0.2224692)
  (8.5,0.2213297)
  (8.75,0.22018356)
  (9,0.21903113)
  (9.25,0.21787295)
  (9.5,0.21670948)
  (9.75,0.21554105)
  (10,0.21436813)
  (10.25,0.21319113)
  (10.5,0.21201037)
  (10.75,0.21082625)
  (11,0.20963914)
  (11.25,0.20844933)
  (11.5,0.20725716)
  (11.75,0.20606297)
  (12,0.20486702)
  (12.25,0.20366962)
  (12.5,0.20247105)
  (12.75,0.20127156)
  (13,0.20007141)
  (13.25,0.19887084)
  (13.5,0.19767009)
  (13.75,0.19646938)
  (14,0.19526891)
  (14.25,0.19406889)
  (14.5,0.19286953)
  (14.75,0.19167099)
  (15,0.19047345)
  (15.25,0.18927709)
  (15.5,0.18808205)
  (15.75,0.18688849)
  (16,0.18569657)
  (16.25,0.1845064)
  (16.5,0.18331811)
  (16.75,0.18213185)
  (17,0.18094769)
  (17.25,0.17976576)
  (17.5,0.17858619)
  (17.75,0.17740902)
  (18,0.17623436)
  (18.25,0.17506234)
  (18.5,0.17389297)
  (18.75,0.17272634)
  (19,0.17156257)
  (19.25,0.17040167)
  (19.5,0.16924369)
  (19.75,0.16808874)
  (20,0.16693683)
  (20.25,0.16578799)
  (20.5,0.16464233)
  (20.75,0.16349982)
  (21,0.1623605)
  (21.25,0.16122447)
  (21.5,0.16009169)
  (21.75,0.15896217)
  (22,0.15783603)
  (22.25,0.15671321)
  (22.5,0.15559371)
  (22.75,0.15447764)
  (23,0.15336493)
  (23.25,0.1522556)
  (23.5,0.15114971)
  (23.75,0.15004721)
  (24,0.14894811)
  (24.25,0.14785247)
  (24.5,0.14676023)
  (24.75,0.14567139)
  (25,0.14458602)
  (25.25,0.14350405)
  (25.5,0.14242547)
  (25.75,0.14135034)
  (26,0.14027862)
  (26.25,0.13921026)
  (26.5,0.13814534)
  (26.75,0.1370838)
  (27,0.13602562)
  (27.25,0.13497084)
  (27.5,0.13391942)
  (27.75,0.13287134)
  (28,0.13182663)
  (28.25,0.13078526)
  (28.5,0.12974719)
  (28.75,0.12871247)
  (29,0.12768106)
  (29.25,0.12665293)
  (29.5,0.12562811)
  (29.75,0.12460658)
  (30,0.1235883)
  (30.25,0.12257331)
  (30.5,0.12156158)
  (30.75,0.12055308)
  (31,0.11954784)
  (31.25,0.11854584)
  (31.5,0.11754706)
  (31.75,0.1165515)
  (32,0.11555918)
  (32.25,0.11457004)
  (32.5,0.11358412)
  (32.75,0.11260142)
  (33,0.1116219)
  (33.25,0.11064557)
  (33.5,0.10967247)
  (33.75,0.10870253)
  (34,0.10773578)
  (34.25,0.10677225)
  (34.5,0.10581189)
  (34.75,0.10485471)
  (35,0.10390076)
  (35.25,0.10294998)
  (35.5,0.10200239)
  (35.75,0.10105803)
  (36,0.10011686)
  (36.25,0.099178883)
  (36.5,0.098244164)
  (36.75,0.097312642)
  (37,0.096384332)
  (37.25,0.095459303)
  (37.5,0.094537493)
  (37.75,0.093618916)
  (38,0.09270365)
  (38.25,0.091791631)
  (38.5,0.09088287)
  (38.75,0.089977458)
  (39,0.089075325)
  (39.25,0.088176485)
  (39.5,0.087281029)
  (39.75,0.086388894)
  (40,0.085500092)
};
\ifhandoutsolutions
\addplot[KSUOrange,dashed,thick,domain=0:40] {0.25-(0.15/35)*x};
\fi
\end{axis}
\end{tikzpicture}
\end{center}
\vspace{-5pt}
By eye, extend the post-xenon trend back to $B=0$. For this exercise, take the line through $(0,0.25)$ and $(35,0.10)$ exactly. Estimate its slope and intercept in Step 5. The curve has no zero crossing in the plotted range.

\Step{4}{State the three LRM assumptions}
\begin{fixedbox}{Assumption 1: identical batches}
Use $N$ equal-mass batches of initially identical fuel; each gains burnup $B_c$ per cycle. For $N=3$, refueling replaces one third of the core.

At \textbf{BOC}, batch burnups are $0,B_c,2B_c$; at \textbf{EOC}, they are $B_c,2B_c,3B_c$. The oldest batch is discharged. In this model, EOC exhausts available excess reactivity.

$B_c$ is the \textbf{cycle burnup}; $B_d=NB_c$ is the \textbf{discharge burnup}. These are burnups, not elapsed times.
\end{fixedbox}

\begin{fixedbox}{Assumption 2: linear fuel reactivity}
\[
 \boxed{\rho_\infty(B)=\rho_0-aB.}
\]
The subscript $\infty$ denotes infinite-medium fuel reactivity. Treat the constant leakage penalty $\rho_L$ separately; $a>0$ is the loss per unit burnup.
\end{fixedbox}

\begin{fixedbox}{Assumption 3: weighted core reactivity}
\[
 \boxed{\rho_{\mathrm{core}}=\sum_{i=1}^{N}w_i\rho_i-\rho_L},
 \qquad \boxed{\sum_{i=1}^{N}w_i=1}.
\]
Here $\rho_i=\rho_\infty(B_i)$; $w_i$ are fixed weights. We use $w_i=1/N$. Adjustable control offsets positive $\rho_{\mathrm{core}}$ during operation.
\end{fixedbox}
\end{minipage}

\newpage
\PageTitle{Batches, cycle length, and discharge burnup}

\noindent\begin{minipage}[t]{0.485\textwidth}\vspace{0pt}\RaggedRight
\Step{5}{Fit a linear fuel-reactivity model}
Use the eyeballed points $(B_1,\rho_1)=(0,0.25)$ and $(B_2,\rho_2)=(35,0.10)$ to approximate the post-xenon trend:
\begin{align*}
 \frac{\Delta\rho}{\Delta B}&=\frac{0.10-0.25}{35-0}
 =\RevealMath[0.65in]{-0.004286},\\
 a&=-\frac{\Delta\rho}{\Delta B}
 =\RevealMath[0.65in]{0.004286},\\
 \rho_0&=\rho_1+aB_1
 =\RevealMath[0.52in]{0.25}.
\end{align*}
Reactivity is dimensionless; $a$ has units $(\mathrm{MWd/kg})^{-1}$. MWd/kg initial HM is numerically equal to GWd/MTU for the initial UO$_2$ fuel. The intercept approximates the post-xenon trend extended back to $B=0$, rather than the fresh-fuel startup value.

\ifhandoutsolutions
\textcolor{HandoutSynthesis}{\textbf{Instructor reference:} Use exactly $\rho_0=0.25$ and $a=0.15/35=3/700$ for this exercise:
$\rho_\infty(B)=0.25-(3/700)B$.
The signed slope is $-0.004285714\ldots$; the intercept is $0.25$. The dashed line shows this graphical approximation. Carry the unrounded ratio $3/700$ into the batch calculations below.}
\par\smallskip
\fi
Use this same two-point approximation for the numerical checks below.

\Step{6}{Test a proposed two-batch cycle}
Take $N=2$, $\rho_L=0.05$, and suppose $B_c=10~\mathrm{GWd/MTU}$.

At BOC, the two batches have burnups $0$ and $B_c$:
\begin{align*}
 \rho_{\mathrm{core}}^{\mathrm{BOC}}
 &=\frac12\left[\rho_\infty(0)+\rho_\infty(B_c)\right]-\rho_L\\
 &=\RevealMath[0.45in]{0.1786}.
\end{align*}
At EOC, the same batches have burnups $B_c$ and $2B_c$:
\begin{align*}
 \rho_{\mathrm{core}}^{\mathrm{EOC}}
 &=\frac12\left[\rho_\infty(B_c)+\rho_\infty(2B_c)\right]-\rho_L\\
 &\simeq\RevealMath[0.45in]{0.1357}.
\end{align*}
Because $\rho_{\mathrm{core}}^{\mathrm{EOC}}>0$, this model permits a longer cycle. Actual cycles also face thermal, fuel-integrity, and operating constraints.

\begin{checkpoint}[Before continuing]
Sketch the two batch burnups at BOC and EOC. Which batch is discharged? Which one is replaced by fresh fuel?
\ifhandoutsolutions
\par\smallskip\textcolor{KSUTeal}{\textbf{Check:} At EOC, the batch at $2B_c$ is discharged and replaced by fresh fuel. The batch at $B_c$ remains for a second cycle.}
\else
\par\vspace{0.40in}
\fi
\end{checkpoint}
\end{minipage}
\hfill
\begin{minipage}[t]{0.485\textwidth}\vspace{0pt}\RaggedRight
\Step{7}{Choose $B_c$ so EOC is critical}
For the two-batch core at EOC,
\begin{align*}
 \bar\rho_\infty^{\mathrm{EOC}}
 &=\frac12\left[(\rho_0-aB_c)+(\rho_0-2aB_c)\right]\\
 &=\rho_0-\frac32aB_c.
\end{align*}
Set $\rho_{\mathrm{core}}^{\mathrm{EOC}}=0$:
\[
 0=\rho_0-\frac32aB_c-\rho_L,
 \qquad
 \boxed{B_c=\frac{2(\rho_0-\rho_L)}{3a}}.
\]
With the eyeballed two-point model,
\begin{align*}
 B_c&\simeq\RevealMath[0.78in]{31.11\ \mathrm{GWd/MTU}},\\
 B_d=2B_c&\simeq\RevealMath[0.78in]{62.22\ \mathrm{GWd/MTU}}.
\end{align*}
For a one-batch core,
\[
 B_c^{(1)}=B_d^{(1)}=\frac{\rho_0-\rho_L}{a}
 \simeq\RevealMath[0.72in]{46.67\ \mathrm{GWd/MTU}}.
\]
\textbf{Range check:} These limiting burnups extrapolate beyond the 40 MWd/kg fit range; they are algebraic LRM results, not recommended fuel-use limits.

Thus, the two-batch core has a shorter cycle,
$B_c^{(2)}=\tfrac23B_c^{(1)}$, but a higher discharge burnup,
$B_d^{(2)}=\tfrac43B_d^{(1)}$--about $33\%$ higher.

\Step{8}{Generalize to an $N$-batch core}
At EOC, the batch reactivities are
\begin{align*}
 \rho_1&=\rho_0-aB_c,\\
 \rho_2&=\rho_0-2aB_c,\\
 &\ \vdots\\[-6pt]
 \rho_N&=\rho_0-NaB_c.
\end{align*}
For equal batch weights,
\begin{align*}
 0&=\rho_{\mathrm{core}}^{\mathrm{EOC}}\\
 &=\frac1N\left[N\rho_0-aB_c(1+2+\cdots+N)\right]-\rho_L.
\end{align*}
Using $1+2+\cdots+N=N(N+1)/2$,
\[
 \boxed{B_c^{(N)}=\frac{2(\rho_0-\rho_L)}{a(N+1)}},
 \qquad
 \boxed{B_d^{(N)}=NB_c^{(N)}}.
\]
For the same $\rho_0$, $a$, and constant $\rho_L$ for all $N$ (with $\rho_0>\rho_L$),
\[
 \boxed{B_c^{(N)}=\frac{2}{N+1}B_c^{(1)}},
 \qquad
 \boxed{B_d^{(N)}=\frac{2N}{N+1}B_d^{(1)}}.
\]
\end{minipage}

\begin{takeawaybox}[What the linear reactivity model says]
Dividing a core into more equal batches makes each individual cycle shorter, but lets the discharged fuel remain in the core for more cycles and reach a larger burnup. The model is intentionally simple: it assumes an equilibrium cycle, equal batch burnup increments, linear fuel reactivity, fixed batch weights, and a constant leakage penalty.
\end{takeawaybox}

\end{document}
